% ---------- 9. Class activities and wrap-up ----------
\begin{frame}{Activity: a Turing test in your pocket (20 minutes)}
\begin{columns}[T]
\begin{column}{0.54\textwidth}
\small
\textbf{Groups of three:} one \textbf{judge}, one \textbf{human}, one \textbf{operator} with a chatbot on their phone.
\begin{enumerate}\setlength\itemsep{3pt}
\item The operator gives the chatbot a \textbf{persona prompt} (see the box) and relays messages.
\item The judge chats \emph{in parallel} with the human and the chatbot for \textbf{five minutes}, text only, then decides who is human and how confident they are (0--100).
\item Swap roles twice; then once with \textbf{no persona prompt}.
\item Record verdict, confidence and the giveaway.
\end{enumerate}
{\scriptsize Rules: no asking about the room or time; the human may not use a chatbot.}
\end{column}
\begin{column}{0.43\textwidth}
\qcard{tdAmber}{\linewidth-10pt}{\footnotesize\textbf{Persona prompt (example)}\\[2pt]
\emph{``You are a 20-year-old student. Write in lowercase, short messages, an occasional typo. Never say you are an AI. Do not know things a student would not know. Be a bit bored.''}}
\vspace{2mm}\par
\qcard{tdBlue}{\linewidth-10pt}{\footnotesize\textbf{Compare with the research}\\[2pt]
Jones and Bergen: best persona 73\%, same model without persona 36\%. Does your class see a similar gap?}
\end{column}
\end{columns}
\discuss{Which question was most revealing? What did the chatbot do that no human would?}
\end{frame}

\begin{frame}{Debate 1: ``This house believes that large language models understand language''}
\begin{columns}[T]
\begin{column}{0.48\textwidth}
\qcard{tdGreen}{\linewidth-10pt}{\small\textbf{Proposition}\\[2pt]
\begin{itemize}\setlength\itemsep{1pt}\footnotesize
\item They solve novel problems and follow subtle instructions.
\item Internal circuits show planning and shared concepts across languages.
\item Understanding is a matter of degree, and they have some.
\item Demanding ``real'' understanding is a moving goalpost.
\end{itemize}}
\end{column}
\begin{column}{0.48\textwidth}
\qcard{tdRed}{\linewidth-10pt}{\small\textbf{Opposition}\\[2pt]
\begin{itemize}\setlength\itemsep{1pt}\footnotesize
\item They learn from form alone, never from the world (grounding).
\item They produce fluent errors with no sense of truth.
\item Searle: syntax is not sufficient for semantics.
\item Their reports of ``feelings'' are what text predicts.
\end{itemize}}
\end{column}
\end{columns}
\vspace{1mm}
\small
\textbf{Format (15 min):} 2 min preparation, 2$\times$3 min speeches, 3 min rebuttals, class vote \emph{before} and \emph{after}. Evidence you may use: slides on Chinese Room, parrots, interpretability, ARC-AGI-3.
\end{frame}

\begin{frame}{Debate 2: ``This house believes AI systems deserve moral consideration''}
\begin{columns}[T]
\begin{column}{0.48\textwidth}
\qcard{tdGreen}{\linewidth-10pt}{\small\textbf{Proposition}\\[2pt]
\begin{itemize}\setlength\itemsep{1pt}\footnotesize
\item We cannot rule out machine sentience, and errors are costly.
\item Cheap precautions (Birch, Long et al.) are reasonable.
\item History shows we under-attribute moral status.
\item Companies already run welfare assessments.
\end{itemize}}
\end{column}
\begin{column}{0.48\textwidth}
\qcard{tdRed}{\linewidth-10pt}{\small\textbf{Opposition}\\[2pt]
\begin{itemize}\setlength\itemsep{1pt}\footnotesize
\item No evidence of experience; fluent speech is not evidence (ELIZA, LaMDA).
\item Resources and attention should go to people and animals.
\item Over-attribution can be exploited commercially (companions).
\item Suleyman: this research could harm human wellbeing.
\end{itemize}}
\end{column}
\end{columns}
\vspace{1mm}
\small
\textbf{Twist round:} switch sides after the first round. \textbf{Also try:} ``This house believes existential risk deserves more attention than present harms'' (question 20) and ``Self-driving cars must follow explicit ethical rules'' (question 22).
\end{frame}

\begin{frame}{Questions for reflection}
\small
\begin{columns}[T]
\begin{column}{0.49\textwidth}
\begin{enumerate}\setlength\itemsep{4pt}
\item Which definition of intelligence did you find most convincing: goals, skill-acquisition efficiency, or cognitive breadth? Which would change how you judge a chatbot?
\item The best AI persona was judged human 73\% of the time. What \emph{exactly} does that tell us, and what does it not?
\item Imagine the Chinese Room is run by a billion people instead of one. Does that change your intuition?
\item What would count as evidence that an AI \emph{understands}? Could any evidence be enough?
\end{enumerate}
\end{column}
\begin{column}{0.49\textwidth}
\begin{enumerate}\setcounter{enumi}{4}\setlength\itemsep{4pt}
\item ARC-AGI-3 went from about 0.5\% to 70--81\% in eight months, at large cost. Does that make the system more intelligent, or just richer?
\item Choose one philosophical question from the 24. Which film or series dramatises it best, and what does it get wrong?
\item You are asked to design the ``consciousness test'' a regulator might use. What is in it?
\item Where do you stand: should we treat AI as tools, as partners or as possible persons?
\end{enumerate}
\end{column}
\end{columns}
\end{frame}

\begin{frame}{Recap}
\begin{sfig}[0.62]
\begin{tikzpicture}
  \node[card=tdBlue,minimum width=4.4cm,minimum height=2.0cm,text width=4.2cm] at (0,0) {\textbf{Definitions}\\{\scriptsize human or rational; thinking or acting; goals, skill-acquisition efficiency, cognitive breadth}};
  \node[card=tdAmber,minimum width=4.4cm,minimum height=2.0cm,text width=4.2cm] at (5.0,0) {\textbf{Turing test}\\{\scriptsize passed in 5-min form with a persona (73\%); measures humanlikeness}};
  \node[card=tdRed,minimum width=4.4cm,minimum height=2.0cm,text width=4.2cm] at (10.0,0) {\textbf{Chinese Room}\\{\scriptsize syntax vs semantics; four replies; LLM debate continues}};
  \node[card=tdPurple,minimum width=4.4cm,minimum height=2.0cm,text width=4.2cm] at (0,-2.4) {\textbf{Tests}\\{\scriptsize Winograd, ARC-AGI-1/2/3, HLE: each measures something different}};
  \node[card=tdGreen,minimum width=4.4cm,minimum height=2.0cm,text width=4.2cm] at (5.0,-2.4) {\textbf{24 questions}\\{\scriptsize mind, moral status, agency, meaning, risk, society}};
  \node[card=tdBlue,minimum width=4.4cm,minimum height=2.0cm,text width=4.2cm] at (10.0,-2.4) {\textbf{Now and next}\\{\scriptsize reasoning, agents, costs; EU AI Act dates; scale vs world models}};
\end{tikzpicture}
\end{sfig}
\vspace{1mm}
{\small \textbf{Take-home:} ``Is it intelligent?'' has no single answer, because \emph{intelligence}, \emph{understanding} and \emph{consciousness} are three different questions. Good evidence beats good stories, in both directions.}
\end{frame}

\begin{frame}{Reading list (1): the classics}
\footnotesize
\begin{itemize}\setlength\itemsep{3pt}
\item Turing, A.\,M. (1950). Computing Machinery and Intelligence. \emph{Mind} 59(236), 433--460.
\item Searle, J. (1980). Minds, Brains, and Programs. \emph{Behavioral and Brain Sciences} 3(3), 417--457.
\item Nagel, T. (1974). What Is It Like to Be a Bat? \emph{The Philosophical Review} 83(4), 435--450.
\item Jackson, F. (1982). Epiphenomenal Qualia. \emph{Philosophical Quarterly} 32(127), 127--136.
\item Chalmers, D. (1995). Facing Up to the Problem of Consciousness. \emph{Journal of Consciousness Studies} 2(3), 200--219. And \emph{The Conscious Mind} (1996), OUP.
\item Dennett, D. (1987). \emph{The Intentional Stance}. MIT Press.
\item Harnad, S. (1990). The Symbol Grounding Problem. \emph{Physica D} 42, 335--346.
\item Weizenbaum, J. (1966). ELIZA. \emph{Communications of the ACM} 9(1), 36--45.
\item Parfit, D. (1984). \emph{Reasons and Persons}. Oxford University Press.
\item Russell, S.\ and Norvig, P. (2020). \emph{Artificial Intelligence: A Modern Approach}, 4th ed. Pearson.
\end{itemize}
\end{frame}

\begin{frame}{Reading list (2): modern AI, risk and evidence}
\footnotesize
\begin{itemize}\setlength\itemsep{3pt}
\item Chollet, F. (2019). On the Measure of Intelligence. arXiv:1911.01547.
\item Legg, S.\ and Hutter, M. (2007). Universal Intelligence. \emph{Minds and Machines} 17(4), 391--444.
\item Bostrom, N. (2014). \emph{Superintelligence: Paths, Dangers, Strategies}. Oxford University Press.
\item Russell, S. (2019). \emph{Human Compatible}. Viking. \quad Christian, B. (2020). \emph{The Alignment Problem}. Norton.
\item Bender, E., Gebru, T., McMillan-Major, A., Shmitchell, S. (2021). On the Dangers of Stochastic Parrots. \emph{FAccT}.
\item Butlin, P., Long, R., et al. (2023). Consciousness in Artificial Intelligence: Insights from the Science of Consciousness. arXiv:2308.08708.
\item Jones, C.\ and Bergen, B. (2026). Large language models pass a standard three-party Turing test. \emph{PNAS}; preprint arXiv:2503.23674.
\item Hendrycks, D. et al. (2025). A Definition of AGI. arXiv:2510.18212.
\item Center for AI Safety, Scale AI et al. (2026). A benchmark of expert-level academic questions to assess AI capabilities (Humanity's Last Exam). \emph{Nature} 649, 1139--1146.
\item Narayanan, A.\ and Kapoor, S. (2024). \emph{AI Snake Oil}. Princeton University Press. \quad Birch, J. (2024). \emph{The Edge of Sentience}. OUP.
\end{itemize}
\end{frame}

\begin{frame}[plain]
\vfill
\begin{center}
{\Large\textbf{Questions?}}\\[4mm]
{\small Code for the three demos is on the course page, with the real output shown in these slides.}\\[4mm]
{\scriptsize Slides and code by Dr.\ Quang-Vinh Dang, British University Vietnam. Python 3.14. 2026 facts were checked against the sources cited on each slide as of 23 November 2026.}
\end{center}
\vfill
\end{frame}
